% s-215-p215-m002; original PDF 215, 216
\Aphorism{s-215-p215-m002}{307}
Die Betonung subjectivistischer Sätze in der Philosophie deutet darauf, dass der Mensch sich im Verhältniss zur sonstigen Natur wieder einmal recht öde, überflüssig, langweilig, recht als Staffage empfindet: wahrscheinlich nicht mit Unrecht. Nun tröstet er sich damit, dass ohne ihn die Berge nicht hoch und die Bäume nicht grün wären; diese mehr als unbestreitbare Wahrheit schmeichelt seiner Eitelkeit. Wäre er Künstler und könnte selber ein Stück Welt schaffen oder darstellen, so hätte er seine Ichheit in einem gegebenen Falle bewiesen und dürfte jenes körperlose, universale, erkenntnisstheoretische Ichgefühl auf sich beruhen lassen.
\par

